\chapter*{Introducción}
\addcontentsline{toc}{chapter}{Introducción}
\markboth{Introducción}{}

El presente documento expone los resultados finales del Trabajo Terminal 2026-B009, desarrollado por Edgar Jair Flores Espinosa, Jose Bryan Omar Jimenez Velazquez y Leilani Shareni Santos Rivera en colaboración con la Defensoría de los Derechos Politécnicos (DDP) del Instituto Politécnico Nacional (IPN). Su propósito fue diseñar e implementar una plataforma web que contribuya a optimizar los procesos de recepción, atención y seguimiento de las quejas que presenta la comunidad politécnica. Para ello, se analizaron las necesidades operativas de la Defensoría, las limitaciones de sus procedimientos y las alternativas tecnológicas que podían contribuir a su mejora. La plataforma integra, además, dos herramientas de inteligencia artificial ---una para la recuperación de antecedentes y la identificación de posibles duplicidades, y otra para la generación de resúmenes narrativos--- como apoyo a las actividades del personal de la Defensoría.

El IPN es una institución pública de educación media superior y superior cuya organización, funcionamiento y atribuciones se encuentran regulados por su marco normativo \cite{ipn_ley_organica}. Dentro de su estructura, la DDP es el órgano encargado de promover, proteger y defender los derechos de la comunidad politécnica \cite{ipn_acuerdo622, ipn_manual_organizacion_ddp}. Para cumplir esta función, recibe y atiende las quejas mediante las cuales estudiantes, personal académico y personal de apoyo dan a conocer hechos que consideran una posible vulneración de sus derechos. Cada queja sigue un procedimiento en el que intervienen sucesivamente las áreas de Recepción, Primer Contacto y Subdefensoría, bajo la supervisión de la persona titular de la Defensoría \cite{ipn_manual_procedimientos_ddp}.

Este procedimiento descansa principalmente en actividades manuales: el registro y la consulta de los expedientes, la lectura completa de cada narrativa, la búsqueda de casos anteriores relacionados con los mismos hechos o personas y la comunicación del avance entre áreas dependen del tiempo y la intervención directa del personal. Ante el volumen de quejas que recibe la Defensoría, esta forma de trabajo incrementa el tiempo de atención de cada caso, provoca la acumulación de expedientes y dificulta que la información necesaria para darles continuidad esté disponible para quien la requiere. A ello se suma que los quejosos no cuentan con un medio para conocer el estado de su queja, lo que en ocasiones los lleva a presentarla de nuevo.

La incorporación de tecnología a este procedimiento no podía limitarse a sustituir los documentos en papel por formularios electrónicos. Debía responder tanto a la gestión administrativa de los expedientes como a las actividades de análisis que realiza el personal, y respetar que las decisiones sobre la competencia, la investigación y la resolución de cada asunto corresponden exclusivamente a la Defensoría. Con ese criterio, la plataforma desarrollada digitaliza el procedimiento completo ---desde la presentación de la queja hasta la conclusión del expediente---, permite al quejoso consultar el estado de su trámite, controla el acceso a la información según el rol de cada persona y hace que el expediente pase completo de un área a otra.

Los dos componentes de inteligencia artificial atienden las actividades de análisis que más tiempo consumen. El primero recupera, a partir de la narrativa de una nueva queja, los expedientes previamente registrados que podrían estar relacionados con ella o ser duplicados, combinando una búsqueda por coincidencia de palabras con una búsqueda por similitud de significado, de modo que puedan encontrarse casos descritos con expresiones distintas. El segundo condensa las narrativas extensas en un resumen breve que conserva los hechos relevantes del relato original. En ambos casos, los resultados son sugerencias que el personal revisa: no sustituyen la lectura del expediente ni la decisión institucional.

La plataforma se construyó con una arquitectura de microservicios, que separa las funcionalidades en componentes independientes, uno por área del procedimiento, y facilita su mantenimiento y despliegue \cite{aws_microservicios, villamizar_microservices}. El desarrollo se apoyó en el levantamiento de requerimientos realizado con el personal de la Defensoría: mediante reuniones de trabajo y la revisión de sus manuales, se modeló el procedimiento en notación BPMN, se definieron los requerimientos y las reglas de negocio de cada actor y se diseñaron las interfaces y la base de datos. La plataforma se implementó de forma incremental, se desplegó en servidores mediante contenedores y se verificó con pruebas de funcionalidad, robustez, fiabilidad, usabilidad y seguridad, así como con la evaluación experimental de sus componentes de inteligencia artificial.

El documento se organiza de la siguiente manera:

\begin{itemize}
	\item El \textbf{Capítulo~\ref{cap:analisis_problema}, Análisis del problema}, describe el dominio de la Defensoría y su procedimiento de atención de quejas, plantea el problema general y los problemas específicos, analiza sus causas y presenta la propuesta de solución.
	\item El \textbf{Capítulo~2, Marco teórico}, presenta el estado del arte de las plataformas de denuncia digital y los fundamentos de procesamiento de lenguaje natural, recuperación de información y generación de resúmenes en los que se basan los componentes de inteligencia artificial.
	\item El \textbf{Capítulo~\ref{cap:proyecto}, Proyecto}, expone los objetivos, la metodología de desarrollo, el alcance de la propuesta, los riesgos y supuestos, los actores, los requerimientos funcionales y no funcionales, las reglas de negocio, los casos de uso, la arquitectura del sistema y el diseño de la base de datos.
	\item El \textbf{Capítulo~\ref{cap:trabajo_realizado}, Trabajo realizado}, muestra las pantallas principales de la plataforma y describe cómo se construyeron sus componentes.
	\item El \textbf{Capítulo~\ref{cap:pruebas}, Pruebas}, documenta la estrategia de pruebas y los casos ejecutados para cada atributo de calidad y para los componentes de inteligencia artificial.
	\item El \textbf{Capítulo~\ref{cap:resultados}, Resultados obtenidos}, resume los defectos encontrados, la evaluación experimental de los componentes de inteligencia artificial y el estado de avance de cada componente.
	\item El \textbf{Capítulo~\ref{cap:futuro_conclusiones}, Trabajo a futuro y conclusiones}, plantea las líneas de continuación del proyecto y valora el cumplimiento de los objetivos.
	\item Finalmente, el \textbf{Anexo~\ref{anx:base_datos}} contiene los diagramas detallados y el diccionario de datos completo de la base de datos, y el \textbf{Anexo~\ref{anx:casos_uso}}, la especificación de todos los casos de uso de la plataforma.
\end{itemize}
